\documentclass[a4paper]{article}

\usepackage[left=1.5cm, right=1.5cm, top=1.5cm, bottom=1.5cm]{geometry} 
\usepackage[utf8]{inputenc} 
\usepackage{german} 
\usepackage{amsmath} 
\usepackage{amssymb} 
\usepackage{listings} 
\usepackage{hyperref} 
\usepackage{graphicx} 
\usepackage{eurosym} 
\usepackage{ulem} \DeclareUnicodeCharacter{20AC}{\euro}

% \usepackage{helvet}
\renewcommand{\familydefault}{\sfdefault}

\begin{document}

\begin{center}
    \begin{minipage}{.2\textwidth}
        \flushleft
        \includegraphics[width=1\linewidth]{Lutherrose.pdf}
    \end{minipage}%
    \begin{minipage}{.6\textwidth}
        \begin{center}
            \footnotesize Verband Christlicher Pfadfinderinnen und Pfadfinder\\
            \large Stamm Martin Luther Lumdatal\\
            ~\\
            \Large \textbf{Änderungsantrag zur Stammesordnung}\\
            \normalsize Anpassung der Ordnung an die gelebte Praxis\\
            ~\\
        \end{center}
    \end{minipage}%
    \begin{minipage}{.2\textwidth}
        \flushright
        \includegraphics[width=.85\linewidth]{Zeichen.pdf}
    \end{minipage}%
\end{center}
~\\[0.5cm] 

\emph{Die Mitgliederversammlung möge beschließen die Stammesordnung wie folgt zu ändern:}

\section{Anpassung der Ordnung in den folgenden Punkten:} % (fold)
\label{sec:anpassung_der_ordnung_auf_die_neue_beauftragung}

\begin{table}[h]
\def\arraystretch{2}
\center
\begin{tabular}{ l|p{.45\textwidth}|p{.45\textwidth}}
\# & \textbf{Stammesordnung 2014}                                                                         & \textbf{geänderte Version}                                                                                                                                                                                                                 \\ \hline
1  & ... Stammeswappen ist nach dem Namensgeber des Stammes die Lutherrose.                               & ... Stammeswappen ist nach dem Namensgeber des Stammes die Lutherrose. \textbf{Sie ist wie in Abbildung 1 zu verwenden.}                                                                                                             \\ \hline
2  & Jede Sippe besitzt \textbf{einen} Sippenführer, \textbf{der} für die Planung und Durchführung ... zuständig \textbf{ist}.       & Jede Sippe besitzt \textbf{mindestens einen} Sippenführer, \textbf{der oder die} für die Planung und Durchführung ... zuständig \textbf{sind}.                                                                                                                        \\ \hline
3  & (nicht vorhanden)                                                                                    & 4.1.1 Aufgaben der Führungsrunde... Betrachtung der langfristigen Planung für den Stamm                                                                                                                                                    \\ \hline
4  & (nicht vorhanden)                                                                                    & 4.1 Führungsrunde... Es ist möglich zwei Personen für eine Beauftragung einzusetzen. Die beiden treten dann als Team gegenüber der MV auf. Sie koordinieren sich selbst und teilen die mit Beauftragung erlangten Aufgaben unter sich auf. \\ \hline
5  & Die Führungsrunde wird vom Stammesführer einberufen und sollte mindestens \textbf{alle 2 Monate} stattfinden. & Die Führungsrunde wird vom Stammesführer einberufen und sollte mindestens \textbf{einmal im Monat} stattfinden.                                                                                                                                    
\end{tabular}
\end{table}
% section anpassung_der_ordnung_auf_die_neue_beauftragung (end)
~\\[0.5cm]
{\emph{Begründung zum Antrag: }
Seit der Beschluss dieser Ordnung auf der ersten Mitgliederversammlung (2011) haben sich einige Kleinigkeiten in der Praxis der Stammesarbeit geändert. Dieser Antrag dient dazu die Ordnung mit der Realität im Stamm abzugleichen. 
\paragraph{zu Punkt 3)} Die langfristige und nachhaltige Planung der Stammesarbeit ist sehr wichtig für die Existenz des Stammes und sollte deshalb als Aufgabe der FüRu gesondert erwähnt werden.


\end{document}






